% !TeX root = ../main.tex
\section{Related Work}
\label{sec:related}

\textbf{Encrypted traffic classification.}
Early work inferred application behavior from flow statistics and hand-crafted
features~\cite{wright2009,anderson2016}. Deep learning replaced manual feature design
with end-to-end models over packet sequences---1D CNNs, recurrent networks, and
attention-based encoders~\cite{wang2017,rezaei2019,traffic_survey,transformer_traffic,vaswani2017}.
QUIC-specific backbone datasets such as CESNET-QUIC22 enable realistic evaluation
rather than synthetic traces~\cite{cesnet_quic22,ring2019}.

\textbf{Traffic morphing and obfuscation.}
Defenses include padding to fixed sizes, adaptive packet scheduling, dummy traffic,
and record-level padding in TLS/QUIC stacks~\cite{rfc9001,dyer2012}. Website
fingerprinting literature formalized attack--defense tradeoffs for morphing
proxies~\cite{walkie_talkie,hayes2016}. Our work focuses on \emph{deterministic},
proxy-realizable transforms---linear block padding, MTU padding, and Laplace timing
jitter---whose overhead can be measured directly from transformed traces.

\textbf{Privacy--utility tradeoffs.}
Prior studies often report accuracy drops without joint analysis of bandwidth/latency
cost~\cite{shusterman2021}. We adopt multi-objective framing and compute
Pareto-efficient operating points, reporting bootstrap confidence intervals and paired
significance tests~\cite{efron1993,mcnemar1947}.

\textbf{Gap.}
Few works jointly (i) evaluate modern QUIC backbone data, (ii) compare Transformer
vs.\ CNN sequence models under identical obfuscation, and (iii) report formal Pareto
frontiers with statistical confidence on held-out temporal test data. We address this
gap with a fully scripted open pipeline.

%==============================================================================
